\documentclass[10pt,aspectratio=1610]{beamer}
%add handout after 10pt, to disable pause
%use aspectratio=1610 to adjust

\usepackage{amssymb}
\usepackage{amsfonts}
\usepackage{amsmath}
\usepackage{amsthm}

\usepackage{mathtools}
\usepackage[T1]{fontenc}
\usepackage{booktabs}
\usepackage{tikz-cd}

\usefonttheme{serif}
\usepackage{newpxtext,newpxmath}

% NOTE: theme name is case-sensitive on Overleaf/Linux ("Berlin", not "berlin").
% The theme must be loaded BEFORE the template settings below, otherwise it overrides them.
\usetheme{Berlin}
\setbeamertemplate{theorems}[numbered]
\setbeamertemplate{navigation symbols}{}

% Operators
\DeclareMathOperator{\Spec}{Spec}
\DeclareMathOperator{\rad}{rad}
\DeclareMathOperator{\Div}{Div}
\DeclareMathOperator{\CaDiv}{CaDiv}
\DeclareMathOperator{\Cl}{Cl}
\DeclareMathOperator{\CaCl}{CaCl}
\DeclareMathOperator{\Pic}{Pic}
\DeclareMathOperator{\divv}{div}
\DeclareMathOperator{\Hom}{Hom}
\DeclareMathOperator{\im}{im}
\DeclareMathOperator{\coker}{coker}
\newcommand{\OO}{\mathcal{O}}
\newcommand{\FF}{\mathcal{F}}
\newcommand{\GG}{\mathcal{G}}
\newcommand{\KK}{\mathcal{K}}
\newcommand{\LL}{\mathcal{L}}
\newcommand{\A}{\mathbb{A}}
\newcommand{\PP}{\mathbb{P}}
\newcommand{\Z}{\mathbb{Z}}
\newcommand{\pp}{\mathfrak{p}}
\newcommand{\mm}{\mathfrak{m}}
\newcommand{\aaa}{\mathfrak{a}}

\title[Divisor of Algebraic Curves\texorpdfstring{\quad\insertframenumber/\inserttotalframenumber}{}]{Divisor of Algebraic Curves}
\subtitle{MA4091 Seminar Matematika I}
\author[G. Allister (10123021)]{Gilbert Allister (10123021) \\
Dosen Pembimbing: Rizal Afgani, Ph.D. \\
Dosen Penguji 1: idk \\
Dosen Penguji 2: idk
}


\date{ \fontsize{8pt}{8pt}\selectfont
  Program Studi Sarjana Matematika \\
  Fakultas Matematika dan Ilmu Pengetahuan Alam \\
  Institut Teknologi Bandung \\[1cm]
  October 5, 2026} % presentation date (was \today)


\begin{document}
    \frame{\titlepage}

    \begin{frame}{Table of Contents}
        \tableofcontents
    \end{frame}

    %=====================================================================
    \section{Introduction}
    %=====================================================================

    \begin{frame}{Motivation}
        \begin{itemize}
            \item Let $k$ be an algebraically closed field. A polynomial $f \in k[x]$ is determined, up to a constant, by its zeros and their multiplicities:
            \begin{align*}
                f = c\,(x-a_1)^{m_1}\cdots(x-a_r)^{m_r}.
            \end{align*}
            \item In the same way, a rational function $f = g/h$ on the projective line $\PP^1 = k \cup \{\infty\}$ is determined, up to a constant, by its \textbf{zeros and poles}. We record them in a formal sum, the \textbf{divisor} of $f$.
            \begin{example}
                For $f = \dfrac{(x-1)^2}{x}$ we get $\divv(f) = 2\,[1] - [0] - [\infty]$: a double zero at $1$ and simple poles at $0$ and $\infty$.
            \end{example}
            \item \textbf{Question.} On an algebraic curve $X$, fix points $P_i$ and integers $n_i$. How large is the space of rational functions with poles of order at most $n_i$ at $P_i$ and no other poles? Divisors are the language for this question; the Riemann--Roch theorem is the answer.
        \end{itemize}
    \end{frame}

    \begin{frame}{Motivation}
        \begin{itemize}
            \item \textbf{Generalizing zeros and poles:} divisors formalize the study of zeros and poles of rational (meromorphic) functions on an algebraic curve.
            \item \textbf{Transition to schemes:} scheme theory extends classical varieties over an algebraically closed field, so that arithmetic situations (arbitrary base rings) and non-reduced structures can be handled in one language.
            \item \textbf{Geometric invariants:} Weil and Cartier divisors are used to construct line bundles, which lead to central results such as the \textbf{Riemann--Roch theorem}.
        \end{itemize}
    \end{frame}

    \begin{frame}{Aim of my Undergraduate Thesis}
        The aims of my undergraduate thesis are:
        \begin{enumerate}
            \item To understand the basic objects of algebraic geometry.
            \item To understand the basics of scheme theory and divisors of algebraic curves.
            \item To identify the role of scheme theory in proving the Grothendieck--Riemann--Roch (GRR) theorem and its corollary, the Riemann--Roch (RR) theorem for curves.
        \end{enumerate}
        \bigskip
        \textbf{This seminar} covers the first two aims: the definitions and the main statements, from affine varieties up to divisors on a curve.
    \end{frame}


    %=====================================================================
    \section{Classical Algebraic Geometry and Affine Scheme}
    %=====================================================================

    \begin{frame}{Outline of this section}
        \begin{enumerate}
            \item Variety
            \item Zariski Topology
            \item Hilbert's Nullstellensatz
            \item Affine Scheme
        \end{enumerate}
    \end{frame}

    \begin{frame}{Variety}
        \begin{definition}[Algebraic set]
            Let $k$ be an algebraically closed field and $A = k[X_1, \ldots, X_n]$. For an ideal $J$ of $A$ and a subset $X \subseteq k^n$, define
            \begin{align*}
                V(J) = \{a \in k^n \mid f(a)=0 \ \forall f \in J\}, \qquad I(X) = \{f \in A \mid f(a) = 0 \ \forall a \in X\}.
            \end{align*}
            $V(J)$ is the \textbf{algebraic set} defined by $J$; $I(X)$ is the \textbf{ideal of} $X$.
        \end{definition}
        \begin{definition}[Affine variety]
            A non-empty algebraic set $X$ is \textbf{irreducible} if it is not the union of two proper algebraic subsets. An \textbf{affine variety} is an irreducible algebraic set; its \textbf{coordinate ring} is $k[X] = A / I(X)$.
        \end{definition}
        \textbf{Example.} In $k^2$: $V(Y - X^2)$ is an affine variety (a parabola); $V(XY) = V(X) \cup V(Y)$ is not irreducible (two lines).
    \end{frame}

    \begin{frame}{Zariski Topology}
        \begin{theorem}
            Let $I$, $J$ and $J_\alpha$ ($\alpha \in \Lambda$) be ideals of $A = k[X_1, \ldots, X_n]$. Then
            \begin{enumerate}
                \item $V(\mathbf{0}) = k^n$ and $V(A) = \emptyset$;
                \item if $I \subseteq J$, then $V(J) \subseteq V(I)$;
                \item $\bigcap_{\alpha \in \Lambda} V(J_\alpha) = V\bigl(\sum_{\alpha \in \Lambda} J_\alpha\bigr)$;
                \item $V(I) \cup V(J) = V(IJ) = V(I \cap J)$.
            \end{enumerate}
        \end{theorem}
        \begin{itemize}
            \item By (1), (3) and (4), the algebraic sets are the closed sets of a topology on $k^n$.
            \begin{definition}[Zariski topology]
                The \textbf{Zariski topology} on $k^n$ is the topology whose closed sets are exactly the algebraic sets $V(J)$.
            \end{definition}
            \item For $n = 1$ the closed sets are $k$ itself and the finite subsets of $k$, so the topology is not Hausdorff.
        \end{itemize}
    \end{frame}

    \begin{frame}{Radical Ideal}
        \begin{itemize}
            \item Before stating the Nullstellensatz, we define the radical of an ideal in order to simplify the statement.
            \begin{definition}[Radical ideal]
                Let $I$ be an ideal of $A = k[X_1, \ldots, X_n]$. The \textbf{radical} of $I$ is
                \begin{align*}
                    \rad I = \{ f \in A : f^m \in I \text{ for some integer } m \geq 1\}.
                \end{align*}
                If $I = \rad I$, we say that $I$ is a \textbf{radical ideal}.
            \end{definition}
            \item $\rad I$ is an ideal containing $I$, and $V(\rad I) = V(I)$.
            \item Every prime ideal is radical, and $I(X)$ is a radical ideal for every $X \subseteq k^n$.
            \begin{example}
                In $k[X]$ we have $\rad\,(X^2) = (X)$. Both ideals define the same algebraic set $\{0\}$, so the algebraic set does not remember the multiplicity.
            \end{example}
        \end{itemize}
    \end{frame}

    \begin{frame}{Hilbert's Nullstellensatz}
        \begin{theorem}[Nullstellensatz, weak form]
            Let $k$ be an algebraically closed field and $A = k[X_1, \ldots, X_n]$. If $\mm$ is a maximal ideal of $A$, then
            \begin{align*}
                \mm = (X_1 - a_1,\, X_2 - a_2,\, \ldots,\, X_n-a_n)
            \end{align*}
            for some point $(a_1, a_2, \ldots, a_n) \in k^n$.
        \end{theorem}

        \begin{theorem}[Nullstellensatz, strong form]
            Let $k$ be an algebraically closed field and $A = k[X_1, \ldots, X_n]$. If $J$ is an ideal of $A$ with $J \neq A$, then $V(J) \neq \emptyset$. Furthermore, for every ideal $J$ of $A$,
            \begin{align*}
                I(V(J)) = \rad J.
            \end{align*}
        \end{theorem}
    \end{frame}

    \begin{frame}{The Algebra--Geometry Dictionary}
        \begin{corollary}
            The maps $X \mapsto I(X)$ and $J \mapsto V(J)$ are mutually inverse, inclusion-reversing bijections:
            \begin{center}
                \upshape
                \begin{tabular}{rcl}
                    \toprule
                    \textbf{Geometry} (in $k^n$) & & \textbf{Algebra} (in $A = k[X_1,\ldots,X_n]$) \\
                    \midrule
                    algebraic sets & $\longleftrightarrow$ & radical ideals \\
                    affine varieties & $\longleftrightarrow$ & prime ideals \\
                    points & $\longleftrightarrow$ & maximal ideals \\
                    \bottomrule
                \end{tabular}
            \end{center}
        \end{corollary}
        \begin{itemize}
            \item So the geometry of $k^n$ can be recovered from the ring $A$ alone.
            \item \textbf{Idea of scheme theory:} replace $A$ by an arbitrary commutative ring with $1$, and take \emph{all} prime ideals as the points of the space.
        \end{itemize}
    \end{frame}

    \begin{frame}{Affine Scheme}
        \begin{definition}[Spectrum of a ring]
            Let $A$ be a commutative ring with $1$. The \textbf{spectrum} of $A$ is the set $\Spec A = \{\pp \subset A \mid \pp \text{ is a prime ideal}\}$. For an ideal $\aaa$ of $A$ and an element $f \in A$, define
            \begin{align*}
                V(\aaa) = \{\pp \in \Spec A \mid \pp \supseteq \aaa\}, \qquad D(f) = \{\pp \in \Spec A \mid f \notin \pp\}.
            \end{align*}
        \end{definition}
        \begin{itemize}
            \item The sets $V(\aaa)$ satisfy the same four properties as before, so they are the closed sets of the \textbf{Zariski topology} on $\Spec A$. The sets $D(f)$ form a basis of open sets.
            \item An \textbf{affine scheme} is $\Spec A$ together with its structure sheaf $\OO_{\Spec A}$, the ``functions'' on $\Spec A$ (made precise in the next section).
        \end{itemize}
        \begin{example}
            $\Spec k[X] = \{(X - a) \mid a \in k\} \cup \{(0)\}$. The points $(X-a)$ are closed; $(0)$ is the \textbf{generic point}, whose closure is the whole space.
        \end{example}
    \end{frame}

    %=====================================================================
    \section{Scheme Theory}
    %=====================================================================

    \begin{frame}{Outline of this section}
        \begin{enumerate}
            \item Presheaf
            \item Sheaf
            \item Morphism of Sheaves
            \item Sheafification
            \item Schemes
            \item Morphism of Schemes
        \end{enumerate}
    \end{frame}

    \begin{frame}{Presheaf}
        \begin{definition}[Presheaf]
            Let $X$ be a topological space. A \textbf{presheaf} $\FF$ of abelian groups on $X$ consists of
            \begin{itemize}
                \item an abelian group $\FF(U)$ for every open set $U \subseteq X$, and
                \item a group homomorphism $\rho_{UV} \colon \FF(U) \to \FF(V)$ for every inclusion $V \subseteq U$ of open sets,
            \end{itemize}
            such that $\FF(\emptyset) = 0$, $\rho_{UU} = \mathrm{id}_{\FF(U)}$, and $\rho_{UW} = \rho_{VW} \circ \rho_{UV}$ whenever $W \subseteq V \subseteq U$.
        \end{definition}
        \begin{itemize}
            \item Elements $s \in \FF(U)$ are called \textbf{sections} of $\FF$ over $U$; we write $s|_V = \rho_{UV}(s)$ and call $\rho_{UV}$ the \textbf{restriction map}. Presheaves of rings are defined in the same way.
            \item The \textbf{stalk} of $\FF$ at a point $p \in X$ is
            \begin{align*}
                \FF_p = \varinjlim_{U \ni p} \FF(U) = \{(U, s) \mid p \in U,\ s \in \FF(U)\} / \sim,
            \end{align*}
            where $(U,s) \sim (V,t)$ if $s|_W = t|_W$ for some open $W$ with $p \in W \subseteq U \cap V$.
        \end{itemize}
    \end{frame}

    \begin{frame}{Sheaf}
        \begin{definition}[Sheaf]
            A presheaf $\FF$ on $X$ is a \textbf{sheaf} if for every open set $U \subseteq X$ and every open cover $\{V_i\}_{i}$ of $U$:
            \begin{enumerate}
                \item (\textbf{Locality}) if $s \in \FF(U)$ satisfies $s|_{V_i} = 0$ for all $i$, then $s = 0$;
                \item (\textbf{Gluing}) if $s_i \in \FF(V_i)$ satisfy $s_i|_{V_i \cap V_j} = s_j|_{V_i \cap V_j}$ for all $i, j$, then there exists $s \in \FF(U)$ with $s|_{V_i} = s_i$ for all $i$.
            \end{enumerate}
        \end{definition}
        By locality, the section $s$ in the gluing axiom is unique.
        \begin{example}
            \begin{itemize}
                \item On any topological space $X$, $\FF(U) = \{\text{continuous functions } U \to \mathbb{R}\}$ with the usual restriction of functions is a sheaf.
                \item For $X = \mathbb{R}$, $\FF(U) = \{\text{bounded continuous functions } U \to \mathbb{R}\}$ is a presheaf but not a sheaf: the functions $s_i(x) = x$ on $V_i = (-i, i)$, $i \geq 1$, are bounded and agree on overlaps, but they glue to an unbounded function on $\mathbb{R}$.
            \end{itemize}
        \end{example}
    \end{frame}

    \begin{frame}[fragile]{Morphism of Sheaves}
        \begin{definition}[Morphism of (pre)sheaves]
            Let $\FF$, $\GG$ be presheaves on $X$. A \textbf{morphism} $\varphi \colon \FF \to \GG$ is a collection of group homomorphisms $\varphi(U) \colon \FF(U) \to \GG(U)$, one for each open set $U \subseteq X$, such that for every inclusion $V \subseteq U$ the following diagram commutes:
            \begin{center}
                \begin{tikzcd}[row sep=1.6em, column sep=3.5em]
                    \FF(U) \arrow[r, "\varphi(U)"] \arrow[d, "\rho_{UV}"'] & \GG(U) \arrow[d, "\rho'_{UV}"] \\
                    \FF(V) \arrow[r, "\varphi(V)"'] & \GG(V)
                \end{tikzcd}
            \end{center}
            A morphism of sheaves is a morphism of the underlying presheaves.
        \end{definition}
        \begin{theorem}
            A morphism $\varphi$ induces homomorphisms on stalks $\varphi_p \colon \FF_p \to \GG_p$. If $\FF$ and $\GG$ are \emph{sheaves}, then $\varphi$ is an isomorphism if and only if $\varphi_p$ is an isomorphism for every $p \in X$.
        \end{theorem}
    \end{frame}

    \begin{frame}[fragile]{Sheafification}
        \begin{itemize}
            \item \textbf{Problem.} For a morphism of sheaves $\varphi \colon \FF \to \GG$, the presheaves $U \mapsto \im \varphi(U)$ and $U \mapsto \coker \varphi(U)$ are in general not sheaves.
        \end{itemize}
        \begin{theorem}[Sheafification]
            Let $\FF$ be a presheaf on $X$. There exist a sheaf $\FF^+$ and a morphism $\theta \colon \FF \to \FF^+$ such that for every sheaf $\GG$ and every morphism $\varphi \colon \FF \to \GG$ there is a unique morphism $\psi \colon \FF^+ \to \GG$ with $\varphi = \psi \circ \theta$:
            \begin{center}
                \begin{tikzcd}[row sep=1.6em, column sep=3.5em]
                    \FF \arrow[r, "\theta"] \arrow[dr, "\varphi"'] & \FF^+ \arrow[d, dashed, "\exists!\, \psi"] \\
                    & \GG
                \end{tikzcd}
            \end{center}
            The pair $(\FF^+, \theta)$ is unique up to unique isomorphism.
        \end{theorem}
        \begin{itemize}
            \item $\FF^+$ is called the \textbf{sheaf associated to} $\FF$. It has the same stalks: $\FF_p \cong \FF^+_p$ for all $p \in X$.
            \item If $\FF$ is already a sheaf, then $\theta$ is an isomorphism.
        \end{itemize}
    \end{frame}

    \begin{frame}{Schemes}
        \begin{definition}[Locally ringed space]
            A \textbf{ringed space} is a pair $(X, \OO_X)$ of a topological space $X$ and a sheaf of rings $\OO_X$ on $X$. It is a \textbf{locally ringed space} if every stalk $\OO_{X,p}$ is a local ring.
        \end{definition}
        \begin{theorem}[Structure sheaf of $\Spec A$]
            Let $A$ be a ring. There is a sheaf of rings $\OO$ on $\Spec A$ such that
            \begin{align*}
                \OO(D(f)) \cong A_f \ \text{ for all } f \in A, \qquad \OO_{\pp} \cong A_{\pp} \ \text{ for all } \pp \in \Spec A.
            \end{align*}
            In particular $\OO(\Spec A) \cong A$, and $(\Spec A, \OO)$ is a locally ringed space.
        \end{theorem}
        \begin{definition}[Scheme]
            An \textbf{affine scheme} is a locally ringed space isomorphic to $(\Spec A, \OO)$ for some ring $A$. A \textbf{scheme} is a locally ringed space $(X, \OO_X)$ in which every point has an open neighbourhood $U$ such that $(U, \OO_X|_U)$ is an affine scheme.
        \end{definition}
    \end{frame}

    \begin{frame}{Examples of Schemes}
        \begin{example}[Affine space]
            $\A^n_k = \Spec k[X_1, \ldots, X_n]$. By the Nullstellensatz, its closed points are the points of $k^n$; its other points are the prime ideals that are not maximal, i.e.\ the generic points of the subvarieties of positive dimension.
        \end{example}
        \begin{example}[Projective line]
            Take two copies of the affine line, $U_0 = \Spec k[x]$ and $U_1 = \Spec k[y]$, and glue the open sets $D(x) \subset U_0$ and $D(y) \subset U_1$ along the isomorphism given by
            \begin{align*}
                k[y, y^{-1}] \xrightarrow{\ \sim\ } k[x, x^{-1}], \qquad y \mapsto x^{-1}.
            \end{align*}
            The result is a scheme $\PP^1_k$. It is not affine: $\OO(\PP^1_k) = k$, while $\PP^1_k$ has more than one point.
        \end{example}
        \begin{itemize}
            \item This is the scheme-theoretic version of $\PP^1 = k \cup \{\infty\}$ from the motivation: $U_0$ contains the finite points and $U_1$ contains $\infty$ (the point $y = 0$).
        \end{itemize}
    \end{frame}

    \begin{frame}{Morphism of Schemes}
        \begin{definition}[Morphism of schemes]
            A \textbf{morphism of schemes} $(X, \OO_X) \to (Y, \OO_Y)$ is a pair $(f, f^\#)$, where
            \begin{itemize}
                \item $f \colon X \to Y$ is a continuous map, and
                \item $f^\# \colon \OO_Y \to f_*\OO_X$ is a morphism of sheaves of rings on $Y$, where $(f_*\OO_X)(V) = \OO_X(f^{-1}(V))$,
            \end{itemize}
            such that for every $p \in X$ the induced map on stalks $f^\#_p \colon \OO_{Y, f(p)} \to \OO_{X,p}$ is a \textbf{local homomorphism}: $(f^\#_p)^{-1}(\mm_{X,p}) = \mm_{Y,f(p)}$.
        \end{definition}
        \begin{theorem}
            Every ring homomorphism $\varphi \colon A \to B$ induces a morphism of schemes $\Spec B \to \Spec A$, $\pp \mapsto \varphi^{-1}(\pp)$, and every morphism of schemes $\Spec B \to \Spec A$ arises from a unique $\varphi$:
            \begin{align*}
                \Hom_{\text{Schemes}}(\Spec B, \Spec A) \cong \Hom_{\text{Rings}}(A, B).
            \end{align*}
        \end{theorem}
        A \textbf{scheme over} $k$ is a scheme $X$ together with a morphism $X \to \Spec k$.
    \end{frame}

    %=====================================================================
    \section{Divisor of Algebraic Curves}
    %=====================================================================

    \begin{frame}{Outline of this section}
        \begin{enumerate}
            \item Algebraic Curves
            \item Weil Divisor
            \item Principal Divisors and the Class Group
            \item Cartier Divisor
            \item Line Bundles and the Riemann--Roch Theorem
        \end{enumerate}
    \end{frame}

    \begin{frame}{Algebraic Curves}
        \begin{definition}[Curve]
            Let $k$ be an algebraically closed field. A \textbf{curve} over $k$ is an integral scheme $X$ of dimension $1$, proper over $k$, all of whose local rings are regular.
        \end{definition}
        \begin{itemize}
            \item In classical language, $X$ is a nonsingular projective curve, for example $\PP^1_k$.
            \item $X$ has a unique generic point $\xi$. The local ring $K(X) = \OO_{X,\xi}$ is a field, the \textbf{function field} of $X$; its elements are the rational functions on $X$.
            \item For every closed point $P \in X$, the local ring $\OO_{X,P}$ is a discrete valuation ring with fraction field $K(X)$. Its valuation
            \begin{align*}
                v_P \colon K(X)^* \to \Z
            \end{align*}
            is the order of vanishing at $P$: $f$ has a \textbf{zero} of order $v_P(f)$ if $v_P(f) > 0$, and a \textbf{pole} of order $-v_P(f)$ if $v_P(f) < 0$.
        \end{itemize}
        \begin{example}
            For $X = \PP^1_k$: $K(X) = k(x)$, and $v_\infty(g/h) = \deg h - \deg g$ for polynomials $g, h \neq 0$.
        \end{example}
    \end{frame}

    \begin{frame}{Weil Divisor}
        Let $X$ be a curve over $k$.
        \begin{definition}[Weil divisor]
            A \textbf{Weil divisor} on $X$ is a formal finite sum
            \begin{align*}
                D = \sum_{P} n_P \, P, \qquad n_P \in \Z,
            \end{align*}
            over the closed points $P$ of $X$, where $n_P = 0$ for all but finitely many $P$. The Weil divisors form the free abelian group $\Div(X)$ generated by the closed points of $X$.
        \end{definition}
        \begin{itemize}
            \item The \textbf{degree} of $D$ is $\deg D = \sum_P n_P \in \Z$. This gives a group homomorphism $\deg \colon \Div(X) \to \Z$.
            \item $D$ is \textbf{effective}, written $D \geq 0$, if $n_P \geq 0$ for all $P$.
        \end{itemize}
        \begin{example}
            On $\PP^1_k$: $D = 2\,[1] - [0] - [\infty]$ is a Weil divisor of degree $0$; it is not effective.
        \end{example}
    \end{frame}

    \begin{frame}{Principal Divisors and the Class Group}
        \begin{definition}[Principal divisor]
            The \textbf{divisor of} a rational function $f \in K(X)^*$ is $\divv(f) = \sum_{P} v_P(f)\, P$; this is a finite sum, since $f$ has only finitely many zeros and poles. A divisor of this form is called \textbf{principal}.
        \end{definition}
        Since $v_P(fg) = v_P(f) + v_P(g)$, the map $\divv \colon K(X)^* \to \Div(X)$ is a group homomorphism.
        \begin{definition}[Linear equivalence, class group]
            Two divisors $D, D'$ are \textbf{linearly equivalent}, written $D \sim D'$, if $D - D' = \divv(f)$ for some $f \in K(X)^*$. The \textbf{divisor class group} is $\Cl(X) = \Div(X) / \{\text{principal divisors}\}$.
        \end{definition}
        \begin{theorem}
            On a curve $X$, every principal divisor has degree $0$: a rational function has as many zeros as poles, counted with multiplicity. Hence $\deg$ induces a homomorphism $\Cl(X) \to \Z$. For $X = \PP^1_k$ it is an isomorphism: $\Cl(\PP^1_k) \cong \Z$.
        \end{theorem}
    \end{frame}

    \begin{frame}{Cartier Divisor}
        \begin{itemize}
            \item Let $\KK$ be the constant sheaf $K(X)$ on $X$, $\KK^*$ its sheaf of invertible elements, and $\OO_X^* \subseteq \KK^*$ the sheaf of units of $\OO_X$. The quotient sheaf $\KK^* / \OO_X^*$ is the sheafification of $U \mapsto \KK^*(U) / \OO_X^*(U)$.
        \end{itemize}
        \begin{definition}[Cartier divisor]
            A \textbf{Cartier divisor} on $X$ is a global section of the sheaf $\KK^* / \OO_X^*$. Concretely, it is given by a family $\{(U_i, f_i)\}_i$, where $\{U_i\}$ is an open cover of $X$ and $f_i \in K(X)^*$, such that
            \begin{align*}
                f_i / f_j \in \OO_X^*(U_i \cap U_j) \quad \text{for all } i, j.
            \end{align*}
            It is \textbf{principal} if it can be given by a single pair $(X, f)$ with $f \in K(X)^*$.
        \end{definition}
        \begin{itemize}
            \item So a Cartier divisor is a divisor that is \emph{locally} the divisor of a rational function; the $f_i$ are its \textbf{local equations}.
            \item Cartier divisors can be defined on any scheme and are directly related to line bundles.
        \end{itemize}
    \end{frame}

    \begin{frame}{Weil Divisors versus Cartier Divisors}
        \begin{theorem}
            Let $X$ be an integral, separated, noetherian scheme all of whose local rings are unique factorization domains. Then the group of Weil divisors is isomorphic to the group of Cartier divisors,
            \begin{align*}
                \Div(X) \cong \CaDiv(X) := \Gamma(X, \KK^*/\OO_X^*),
            \end{align*}
            and principal Weil divisors correspond to principal Cartier divisors.
        \end{theorem}
        \begin{itemize}
            \item This applies to a curve $X$: every $\OO_{X,P}$ is a discrete valuation ring, hence a unique factorization domain. The correspondence is explicit:
            \begin{align*}
                \{(U_i, f_i)\}_i \ \longmapsto\ \sum_{P} v_P(f_{i})\, P \quad (\text{for any } i \text{ with } P \in U_i).
            \end{align*}
            This is well defined: if $P \in U_i \cap U_j$, then $f_i/f_j$ is a unit at $P$, so $v_P(f_i) = v_P(f_j)$.
            \item Consequently $\Cl(X) \cong \CaCl(X) := \CaDiv(X)/\{\text{principal Cartier divisors}\}$.
        \end{itemize}
    \end{frame}

    \begin{frame}{Divisors and Line Bundles}
        \begin{itemize}
            \item To a divisor $D$ with local equations $\{(U_i, f_i)\}$ we associate the subsheaf $\OO_X(D) \subseteq \KK$ generated by $f_i^{-1}$ on $U_i$. It is locally isomorphic to $\OO_X$, i.e.\ an \textbf{invertible sheaf} (line bundle).
            \item $D \sim D'$ if and only if $\OO_X(D) \cong \OO_X(D')$, and $D \mapsto \OO_X(D)$ induces an isomorphism
            \begin{align*}
                \Cl(X) \cong \CaCl(X) \cong \Pic(X),
            \end{align*}
            where $\Pic(X)$ is the group of invertible sheaves up to isomorphism, under $\otimes$.
            \item Global sections answer the question from the motivation:
            \begin{align*}
                L(D) := \Gamma(X, \OO_X(D)) = \{f \in K(X)^* \mid \divv(f) + D \geq 0\} \cup \{0\}, \qquad \ell(D) := \dim_k L(D).
            \end{align*}
        \end{itemize}
    \end{frame}

    \begin{frame}{The Riemann--Roch Theorem}
        \begin{theorem}[Riemann--Roch]
            Let $X$ be a curve of genus $g$ with canonical divisor $K_X$. For every divisor $D$ on $X$,
            \begin{align*}
                \ell(D) - \ell(K_X - D) = \deg D + 1 - g.
            \end{align*}
        \end{theorem}
        \begin{example}
            On $\PP^1_k$ we have $g = 0$. For $n \geq 0$, $L(n[\infty])$ is the space of polynomials in $x$ of degree at most $n$, so $\ell(n[\infty]) = n + 1 = \deg(n[\infty]) + 1 - g$.
        \end{example}
        \begin{itemize}
            \item The theorem relates the number of linearly independent rational functions with prescribed poles to two numbers: the degree of $D$ and the genus of $X$.
            \item Its proof, as a corollary of the Grothendieck--Riemann--Roch theorem, is the goal of the next seminar.
        \end{itemize}
    \end{frame}


    %=====================================================================
    \section{Conclusion and Next Plan}
    %=====================================================================

    \begin{frame}{Conclusion}
        \begin{itemize}
            \item The Nullstellensatz translates the geometry of algebraic sets in $k^n$ into the algebra of ideals of $k[X_1, \ldots, X_n]$. This motivates $\Spec A$, whose points are the prime ideals of an arbitrary ring $A$.
            \item Sheaves describe data that are defined locally and can be glued. A scheme is a locally ringed space that is locally an affine scheme $(\Spec A, \OO)$.
            \item On a curve $X$, every closed point $P$ gives a valuation $v_P$ on the function field. A Weil divisor is a formal sum of closed points, and $\divv(f)$ records the zeros and poles of $f$.
            \item A Cartier divisor is given locally by one rational function. On a curve, Weil and Cartier divisors agree, and
            \begin{align*}
                \Cl(X) \cong \CaCl(X) \cong \Pic(X).
            \end{align*}
            \item The space $L(D)$ of functions with poles bounded by $D$ is the space of global sections of the line bundle $\OO_X(D)$; its dimension is governed by the Riemann--Roch theorem.
        \end{itemize}
    \end{frame}

    \begin{frame}{Next presentation's plan}
        \begin{enumerate}
            \item To apply scheme theory to divisors of algebraic curves in more depth (invertible sheaves, the canonical divisor, cohomology of sheaves).
            \item To prove the Grothendieck--Riemann--Roch (GRR) theorem using scheme theory.
            \item To prove the Riemann--Roch (RR) theorem for curves as a corollary of the GRR theorem.
        \end{enumerate}
    \end{frame}

    \begin{frame}{References}
        \begin{thebibliography}{9}
            \setbeamertemplate{bibliography item}[text]
            \bibitem{hartshorne} R.~Hartshorne, \emph{Algebraic Geometry}, Graduate Texts in Mathematics 52, Springer, 1977.
            \bibitem{liu} Q.~Liu, \emph{Algebraic Geometry and Arithmetic Curves}, Oxford Graduate Texts in Mathematics 6, Oxford University Press, 2002.
            \bibitem{am} M.~F.~Atiyah and I.~G.~Macdonald, \emph{Introduction to Commutative Algebra}, Addison--Wesley, 1969.
            \bibitem{fulton} W.~Fulton, \emph{Algebraic Curves: An Introduction to Algebraic Geometry}, W.~A.~Benjamin, 1969.
        \end{thebibliography}
    \end{frame}

    \begin{frame}
        \centering
        {\Large Thank you}
    \end{frame}

\end{document}